\documentclass[11pt,a4paper]{article}
\usepackage[T1]{fontenc}
\usepackage{lmodern}
\usepackage[margin=29mm]{geometry}
\usepackage{amsmath,amssymb,amsthm,mathtools}
\usepackage{microtype}
\usepackage{enumitem}
\usepackage[colorlinks=true,linkcolor=blue,citecolor=blue,urlcolor=blue]{hyperref}
\hypersetup{pdftitle={The E x CM component of Zilber--Pink for curves in A2},pdfauthor={OpenAI}}
\urlstyle{same}
\Urlmuskip=0mu plus 1mu
\setlist[itemize]{topsep=4pt,itemsep=2pt}
\setlist[enumerate]{topsep=4pt,itemsep=2pt}
\numberwithin{equation}{section}
\newtheorem{theorem}{Theorem}[section]
\newtheorem{proposition}[theorem]{Proposition}
\newtheorem{lemma}[theorem]{Lemma}
\newtheorem{corollary}[theorem]{Corollary}
\theoremstyle{definition}
\newtheorem{definition}[theorem]{Definition}
\theoremstyle{remark}
\newtheorem{remark}[theorem]{Remark}
\newcommand{\A}{\mathcal A}
\newcommand{\Q}{\mathbb Q}
\newcommand{\Z}{\mathbb Z}
\newcommand{\R}{\mathbb R}
\newcommand{\C}{\mathbb C}
\newcommand{\Qbar}{\overline{\mathbb Q}}
\newcommand{\CM}{\mathrm{CM}}
\newcommand{\End}{\operatorname{End}}
\newcommand{\Hom}{\operatorname{Hom}}
\newcommand{\Lie}{\operatorname{Lie}}
\newcommand{\im}{\operatorname{im}}
\newcommand{\disc}{\operatorname{disc}}
\newcommand{\Gal}{\operatorname{Gal}}
\newcommand{\ord}{\operatorname{ord}}
\newcommand{\hF}{h_{\mathrm F}}
\newcommand{\hFp}{h_{\mathrm F}^{+}}
\newcommand{\cplx}{\mathfrak c}
\newcommand{\abs}[1]{\left|#1\right|}
\newcommand{\norm}[1]{\left\|#1\right\|}
\newcommand{\ceil}[1]{\left\lceil#1\right\rceil}

\title{The \(E\times\CM\) component of Zilber--Pink\\
for curves in \(\A_2\)}
\author{OpenAI}
\date{September 24, 2026}

\begin{document}
\maketitle

\begin{abstract}
We prove the \(E\times\CM\) component of Zilber--Pink for Hodge-generic
algebraic curves in \(\A_2\) over \(\Qbar\). Each such curve contains only
finitely many points whose abelian surface is isogenous to a product of
elliptic curves with at least one factor having complex multiplication.
\end{abstract}

\section{Introduction}

Let \(\A_2\) be the coarse moduli variety of principally polarized abelian
surfaces. Its dimension is three. The Zilber--Pink conjecture predicts
that a Hodge-generic curve has only finitely many intersections with the
union of special subvarieties of codimension at least two. Among these
special subvarieties are the \(E\times\CM\) curves: one elliptic factor
varies, the other has complex multiplication, and one allows all
polarized Hecke translates. Here and below, special subvarieties are
Shimura special subvarieties, including their Hecke translates.
For \(s\in\A_2(\C)\), write \(A_s\) for a principally polarized
abelian surface representing its moduli point. An elliptic curve
has CM if its geometric endomorphism algebra is an imaginary
quadratic field.
In a threefold, a curve and a codimension-two subvariety have
negative expected intersection dimension. The conjecture controls
the exceptional intersections even as the special subvariety varies.
It belongs to the unlikely-intersection framework initiated by
Zilber's atypical-intersection conjectures in semiabelian varieties
\cite{Zilber} and formulated for mixed Shimura varieties by Pink
\cite[Conjecture~1.3]{Pink}.

\begin{theorem}\label{thm:main}
Let \(C\subset\A_2\) be a reduced, irreducible, closed algebraic curve
defined over \(\Qbar\). Suppose that \(C\) is Hodge generic, meaning that
it is contained in no proper special subvariety of \(\A_2\).
Then the set
\[
 \bigl\{s\in C(\Qbar):
 A_s\text{ is isogenous to }E_1\times E_2
 \text{ with at least one }E_i\text{ having CM}\bigr\}
\]
is finite.
\end{theorem}

The isogenies in Theorem~\ref{thm:main} are
isogenies of abelian varieties; they are not required to respect the
given polarization. All isogeny degrees, imaginary quadratic fields,
and endomorphism orders are included, as is the case of two CM factors.
There is no assumption on the boundary of \(C\), on an integral base
point, or on reduction at any finite prime. Thus
Theorem~\ref{thm:main} resolves positively the \(E\times\CM\) component
of Zilber--Pink for Hodge-generic curves in \(\A_2\).

Daw and Orr proved this finiteness statement for curves whose closure
meets the zero-dimensional stratum of the Baily--Borel boundary
\cite[Theorem~1.1]{DO}. They also proved that, without this boundary
hypothesis, a suitable lower bound for Galois orbits suffices
\cite[Theorem~1.2]{DO}. Their later work proves full Zilber--Pink
for curves in \(\A_2\) under the same zero-dimensional-boundary
hypothesis \cite[Corollary~1.2]{DOMultiplicative}.
Papas obtains height and finiteness results for Hodge-generic families
having an interior fiber isogenous to a square of a non-CM elliptic curve
or to a product of one CM and one non-CM elliptic curve. That fiber must have
potentially good reduction at every finite place, and the finiteness
statement imposes a uniform bound on the number of supersingular
places of proximity \cite[Theorem~1.5 and Corollary~1.6]{Papas}.
The both-CM case follows
from Andr\'e--Oort for \(\A_2\), proved by Pila and Tsimerman
\cite[Theorem~1.1]{PT} and included in Tsimerman's general theorem
for \(\A_g\) \cite[Theorem~5.3]{Tsimerman}.
Our task is to supply the orbit bound
for the entire mixed locus. We call a point \emph{mixed} when its
abelian surface is isogenous to a product of one CM and one non-CM
elliptic curve.

The passage from Galois orbits to finiteness is an instance of the
strategy of Pila and Zannier \cite{PZ}: compare arithmetic lower bounds
for Galois orbits with the subpolynomial bounds in the Pila--Wilkie
theorem \cite{PW}. Daw--Orr establish the required geometric and
counting implication for this locus. The arithmetic estimates here
use the theory of periods and minimal abelian subvarieties developed
by Masser and W\"ustholz \cite{MW}, in the arbitrary-polarization
form of Gaudron and R\'emond \cite{GR}. The degeneration argument
uses the integral compactifications of Faltings--Chai and Lan
\cite{FC,Lan}, within the reduction and monodromy framework of
Grothendieck \cite{SGA7} and Tate's explicit elliptic uniformization
\cite{Tate}.

There are two arithmetic ingredients. First, the product-isogeny
description of Daw--Orr \cite[Lemmas~2.1 and~3.2]{DO} gives a mixed
principally polarized surface a canonical graph presentation
\[
 E\times F\longrightarrow A,\qquad \deg(E\times F\longrightarrow A)=n^2,
\]
where \(E\) is non-CM, \(F\) is CM, and the pullback polarization is
\(n\) times the product polarization. At a finite place where \(E\) is
a Tate curve, the rank-one degeneration length of \(A\) is
\(\log|j_E|/n\). Combining this reciprocal factor with the height
machine on a fixed curve gives
\begin{equation}\label{eq:intro-height}
 h(j_E)\ll_C n(1+\sqrt D),
 \qquad D=\abs{\disc\End(F)},
\end{equation}
for all sufficiently large \(n\). The argument works both when the
curve meets the boundary and when it does not.

Second, we use the gluing data to construct an abelian threefold
isogenous to \(E^2\times F\), with an essential period involving both
periods of \(E\). A weighted polarization balances their possibly very
different lengths. The period theorem of Gaudron and R\'emond then
produces a factor \(n^{4/3}\), which exceeds the linear dependence on
\(n\) in \eqref{eq:intro-height}. This construction and the uniform
height estimate are the additional inputs beyond the existing
unlikely-intersection theorem.

We obtain the following explicit, deliberately coarse form of the
arithmetic conclusion. At a mixed point \(s\), let \(N(s)\) be the
minimum degree of an isogeny from an elliptic product to \(A_s\).
The CM factor in a product attaining this minimum is unique up to
isomorphism; this will be proved in Lemma~\ref{lem:gluing}. Write
\[
 \cplx(s)=\max\{N(s),\,|\disc\End(F)|\}.
\]

\begin{theorem}\label{thm:orbits}
Let \(C\) be as in Theorem~\ref{thm:main}, and let \(K\) be a number
field over which it is defined. There is a constant \(c>0\), depending
on \(C\) and \(K\), such that every mixed point \(s\in C(\Qbar)\)
satisfies
\[
 \#\bigl(\Gal(\Qbar/K)\cdot s\bigr)
 \ge c\,\cplx(s)^{1/42}.
\]
\end{theorem}

The proof below uses Siegel's potentially ineffective class-number
bound. Remark~\ref{rem:effective-cm} gives an effective alternative for
that input alone; we make no effectivity claim for the constants in
Theorem~\ref{thm:orbits}. The exponent records the positive power needed
by the counting theorem; its optimization is not needed here.
Section~\ref{sec:tools} fixes the height and period normalizations and
collects the CM estimates. Section~\ref{sec:gluing} proves the canonical
gluing description and its reciprocal degeneration formula.
Section~\ref{sec:height} turns that formula into the height bound on
the fixed curve, treating an empty boundary separately.
Section~\ref{sec:threefold} constructs and balances the essential
period. Section~\ref{sec:orbits} combines the estimates to prove
Theorem~\ref{thm:orbits}, and Section~\ref{sec:finiteness} applies
the corrected Daw--Orr implication and Andr\'e--Oort.

\section{Heights, periods, and CM elliptic curves}\label{sec:tools}

\subsection{Conventions}

Endomorphism rings and homomorphism groups are geometric unless a
ground field is indicated. All heights are absolute logarithmic heights. We write \(h(j_P)\)
for the usual height of the \(j\)-invariant of an elliptic curve \(P\),
and \(\hF(B)\) for stable Faltings height. Set
\(\hFp(B)=\max\{1,\hF(B)\}\).
For a line bundle \(L\) on a \(g\)-dimensional abelian variety,
\(\deg_L B=c_1(L)^g[B]\). Thus a principal polarization has degree \(g!\)
in this convention. The associated period norm is the norm of its
positive Hermitian form.

For an elliptic curve with its principal polarization, choose a period
basis \(e,e\tau\) with \(\tau\) in the standard closed fundamental region.
Writing \(y=\Im\tau\), one has
\begin{equation}\label{eq:elliptic-norms}
 y\ge\sqrt3/2,\qquad
 \|e\|^2=1/y,\qquad
 \|e\tau\|^2=|\tau|^2/y\ll y.
\end{equation}
The Fourier expansion of \(j\), together with compactness away from
the cusp, gives
\begin{equation}\label{eq:j-y}
 \abs{\log^+|j(\tau)|-2\pi y}\ll1.
\end{equation}
Constants in these two formulas are absolute.
Other implicit constants may depend on a fixed curve, family, and
ground field, but never on a varying point or its field of definition.

For a number field \(k\), absolute values at finite places extend the
standard absolute values on \(\Q\); they are not renormalized after a
field extension. Heights are sums with weights
\([k_v:\Q_v]/[k:\Q]\). Equivalently, the archimedean contribution is
the average over all complex embeddings, with conjugate embeddings
counted separately.

\subsection{Two established height and period estimates}

We use product additivity and the isogeny inequality for stable
Faltings height:
\begin{equation}\label{eq:faltings-isogeny}
 \hF(B_1\times B_2)=\hF(B_1)+\hF(B_2),\qquad
 |\hF(B_1)-\hF(B_2)|\le\tfrac12\log\deg\phi
\end{equation}
for an isogeny \(\phi:B_1\to B_2\); see \cite[\S2.3]{GR}.
For elliptic curves, \(\hF\) is bounded below and
\begin{equation}\label{eq:elliptic-height}
 1+h(j_P)\asymp 1+\hFp(P).
\end{equation}
These are stable-height statements, so passing to a semistable
extension does not change the heights. For the comparison with
\(j\)-height, see Silverman \cite[Proposition~2.1]{Silverman}
and the normalization recorded in \cite[Proposition~3.1]{Lobrich}.

\begin{lemma}\label{lem:family-height}
Let \(\mathcal X\to S\) be a fixed principally polarized abelian scheme
over a smooth curve defined over a number field, and let \(\bar S\)
be its smooth projective completion. If \(H\) is an ample divisor on
\(\bar S\), and \(h_H\) is shifted to be nonnegative, then
\[
 \hFp(\mathcal X_t)\ll 1+h_H(t)
 \qquad(t\in S(\Qbar)).
\]
\end{lemma}

\begin{proof}
After a fixed finite cover and restriction to a dense open, choose
a symmetric ample line bundle representing the polarization and
theta data as in \cite[\S2.3]{Pazuki}, for one fixed even integer
\(r\ge2\). This is possible by choosing the data over a finite
extension of the function field and spreading out. The theta-null map is an algebraic map to projective space.
Pazuki's theta/Faltings comparison
\cite[Theorem~1.1 and Corollary~1.3]{Pazuki} gives
\(\hFp\ll1+h_\Theta\), with constants depending only on the dimension
and the fixed theta level.

The theta-null map extends to the smooth projective covering curve.
If \(p\) is the covering map and \(T\) is the pullback of the
hyperplane divisor, choose an integer \(a\) such that
\(ap^*H-T\) is ample. Its height is bounded below, so the height
machine gives \(h_\Theta\le a h_H\circ p+O(1)\).
This proves the asserted bound after lifting \(t\); stable Faltings
height is unchanged by the lift. The finite set omitted when choosing the data can be included
by enlarging the constant.
\end{proof}

Call a nonzero period \(\omega\) of a complex abelian variety
\emph{essential} if it belongs to the Lie algebra of no proper
abelian subvariety. For a polarized abelian variety \((B,L)\) of
dimension \(g\) over a number field \(k\), and an essential period
\(\omega\) at one complex
embedding of \(k\), \cite[Theorem~1.2]{GR} gives
\begin{equation}\label{eq:period-theorem}
 \frac{(\deg_L B)^{1/g}}{\|\omega\|_L^2}
 \ll_g [k:\Q]\max\{1,\hF(B),\log\deg_L B\}.
\end{equation}
To explain the field-degree factor, let
\(\delta_0(B_\sigma,L_\sigma)\) be the least norm of an essential
period at the embedding \(\sigma\), with value \(+\infty\) if
there is none. The essential-minimum form of their theorem bounds
\[
 \frac1{[k:\Q]}\sum_{\sigma:k\hookrightarrow\C}
 \frac{(\deg_LB)^{1/g}}{\delta_0(B_\sigma,L_\sigma)^2}
 \ll_g \max\{1,\hF(B),\log\deg_LB\}.
\]
At the chosen embedding,
\(\delta_0(B_\sigma,L_\sigma)\le\|\omega\|_L\).
Keeping this one nonnegative summand proves
\eqref{eq:period-theorem}. The estimate allows arbitrary
polarization degree.

\subsection{CM estimates}\label{subsec:cm}

\begin{lemma}\label{lem:cm}
Let \(k\) be a number field and \(F/k\) a CM elliptic curve. Set
\(\mathcal O=\End_{\Qbar}(F)\), \(D=|\disc\mathcal O|\), and \(d=[k:\Q]\).
Then
\begin{equation}\label{eq:cm-bounds}
 D\ll d^4,\qquad h(j_F)\ll1+\sqrt D.
\end{equation}
At any complex embedding, choose a reduced period basis \(f,f\tau_F\)
and put \(y_F=\Im\tau_F\). Then
\begin{equation}\label{eq:b-index}
 b=[\Lambda_F:\mathcal O f]
   =\frac{\sqrt D}{2y_F}\in\Z_{>0},\qquad
 y_F\le\sqrt D/2,\qquad b\ll\sqrt D.
\end{equation}
\end{lemma}

\begin{proof}
The degree of \(j_F\) over \(\Q\) is the class number \(h(\mathcal O)\),
and \(j_F\) is an algebraic integer \cite[Theorem~11.1 and Proposition~13.2]{Cox}.
Write \(D=f_0^2D_0\), where \(D_0\) is the absolute discriminant of
the maximal order. Siegel's class-number bound gives
\(h(\mathcal O_{D_0})\gg D_0^{1/4}\) \cite{Siegel}.
The conductor formula \cite[Theorem~7.24]{Cox} multiplies this by
\[
 \frac{f_0}{[\mathcal O_{D_0}^{\times}:\mathcal O^\times]}
 \prod_{p\mid f_0}\left(1-\frac{\chi_{-D_0}(p)}p\right),
\]
which is at least \(\varphi(f_0)/3\gg\sqrt{f_0}\).
Here \(\chi_{-D_0}\) is the quadratic character of the imaginary
quadratic field of discriminant \(-D_0\), defined by the Kronecker
symbol, and \(\varphi\) is Euler's totient function. The elementary bound
\(\varphi(f_0)\ge\sqrt{f_0/2}\) suffices.
Thus \(d\ge h(\mathcal O)\gg D^{1/4}\), including every conductor.
The constant may be ineffective.

The covolume of \(\mathcal O\) in \(\C\) is \(\sqrt D/2\), whereas
that of \(\Lambda_F/f=\Z+\tau_F\Z\) is \(y_F\). Since
\(\mathcal O f\subset\Lambda_F\), their ratio gives
\eqref{eq:b-index}. The lower bound for \(y_F\) in
\eqref{eq:elliptic-norms} gives \(b\ll\sqrt D\).
Finally, \eqref{eq:j-y} bounds every archimedean conjugate of \(j_F\)
by \(O(1+\sqrt D)\); integrality removes all finite-place contributions.
\end{proof}

\begin{remark}[An effective alternative for the CM bound]
\label{rem:effective-cm}
The Dirichlet assertion of \cite[Theorem~1.1]{OpenAIQuasiRH2026}
makes the implied constant in \(D\ll d^4\) effective. Indeed, its
zero-free half-plane \(\operatorname{Re}s>7/8\) supplies the short real
zero-free interval used in
\cite[proof of Theorem~2, pp.~363--364]{Chen2007}. The resulting estimate is
\(L(1,\chi_{-D_0})\gg1/\log(2D_0)\), with an effective absolute
constant, after handling finitely many small conductors. The Dirichlet
class-number formula therefore gives
\(h(\mathcal O_{D_0})\gg\sqrt{D_0}/\log(2D_0)\gg D_0^{1/4}\)
effectively. The conductor estimate in the preceding proof then yields
\(d\ge h(\mathcal O)\gg D^{1/4}\) with an effective constant for
every order. This concerns only the CM discriminant estimate: it does
not assert effective Galois-orbit constants or effective finiteness in
Theorems~\ref{thm:orbits} and~\ref{thm:main}.
\end{remark}

\section{Canonical gluing and boundary lengths}\label{sec:gluing}

\subsection{The canonical product isogeny}

We first identify the product and gluing integer intrinsically, so
that the subsequent height bounds use the same complexity as the
Galois-orbit theorem. The description below follows
\cite[Lemmas~2.1 and~3.2]{DO}; we include the argument to make the
minimality and field of definition explicit.

\begin{lemma}\label{lem:gluing}
Let \((A,\lambda)\) be a principally polarized abelian surface over
a field of characteristic zero, and suppose that \(A\) is isogenous
geometrically to the product of a non-CM and a CM elliptic curve.
There are unique elliptic subvarieties \(E,F\subset A\), with \(E\)
non-CM and \(F\) CM. They are defined over the field of \(A\).
For some positive integer \(n\), addition induces an isogeny
\[
 \phi:E\times F\longrightarrow A
\]
of degree \(n^2\), whose pullback polarization is \(n\) times the
product polarization and whose kernel is the graph of an isomorphism
\(\theta:F[n]\to E[n]\). This isomorphism is defined over the field of
\(A\). The minimum degree of an isogeny from an elliptic product to
\(A\) is \(n^2\), and its factors are then isomorphic to \(E,F\).
\end{lemma}

\begin{proof}
The two simple isogeny factors are nonisogenous, so there are no
homomorphisms between them. Their images in \(A\) are therefore
unique elliptic subvarieties. Uniqueness gives descent.
The pullback of \(\lambda\) to \(E\times F\) is diagonal, say with
positive elliptic multipliers \(n_E,n_F\). Comparing degrees of
polarizations gives \(|\ker\phi|=n_En_F\).
Each projection of the kernel is injective, since \(E,F\) are
embedded subvarieties, and
\[
 \ker\phi\subset E[n_E]\times F[n_F].
 \]
Consequently \(n_En_F\le n_E^2\) and \(n_En_F\le n_F^2\), which force
\(n_E=n_F=n\). The projections then identify the kernel with each
full \(n\)-torsion group. This gives the graph, including its descent.

Any isogeny \(P\times Q\to A\) from an elliptic product maps the two
factors, in some order, onto \(E,F\). It factors through their
addition isogeny, so its degree is at least \(n^2\). Equality forces
both factor maps to have degree one.
\end{proof}

The Galois action on the endomorphism ring of a CM elliptic curve
has image of order at most two: it acts on its imaginary quadratic
endomorphism algebra. Thus, when \(A\) is defined over a number field
\(k_1\), all endomorphisms of \(F\) are defined after an extension of
\(k_1\) of degree at most two. Defining the graph itself requires
no such extension.

\subsection{The degeneration input}\label{subsec:degeneration}

The new local issue is how the integer \(n\) changes the degeneration
of the non-CM factor. We first compute the resulting length from
semiabelian uniformization, then compare it with the boundary of a
fixed integral compactification.

Suppose a principally polarized abelian variety over a complete
discretely valued field has split semiabelian reduction. Its
uniformization consists of an extension of a good-reduction abelian
variety by a split torus, modulo a period lattice \(Y\).
If \(X\) is the character lattice of the torus, a polarization gives
a map \(Y\to X\), an isomorphism for a principal polarization.
The degeneration data carry a functorial valuation pairing
\[
 Y\times X\longrightarrow\R.
\]
We use the extended absolute value in defining this pairing. In toric
rank one, its value on generators identified by the principal
polarization is a positive number, denoted \(\ell(A)\).

\begin{lemma}\label{lem:length}
In the situation of Lemma~\ref{lem:gluing} over a number field, let
\(v\) be a finite place. If \(E\) has potential good reduction, so
does \(A\). Otherwise, after an extension giving split semiabelian
reduction,
\begin{equation}\label{eq:length}
 \ell(A)=\frac{\log|j_E|_v}{n}.
\end{equation}
\end{lemma}

\begin{proof}
The CM factor has potential good reduction
\cite[Theorem~6(a)]{ST}; for an elliptic curve this also follows from
the integrality of its \(j\)-invariant. After a finite local extension,
it has good reduction and \(E\) has either good reduction or split
Tate reduction \cite[Theorem~1 and Applications]{Tate}.
Potential good reduction and rational toric rank are isogeny
invariants. It remains to determine the length in the Tate case.

Functoriality of semiabelian uniformization gives maps of tori and
period lattices for \(\phi:E\times F\to A\)
\cite[Corollary~4.5.5.5]{Lan}. The map of one-dimensional tori is an
isomorphism. Indeed, a nontrivial kernel would contain nontrivial
roots of unity. Torus torsion injects into the Tate elliptic curve,
because its period lattice is torsion free. Such a kernel would
therefore give a nonzero point of \(\ker\phi\) in \(E\times\{0\}\),
contradicting the graph description.
This reasoning takes place on characteristic-zero generic fibers
and does not require \(v(n)=0\).

Write \(Y',X'\) for the source lattices of \(E\times F\), with
its product principal polarization. Pulling back the polarization
means that the composite
\[
 Y'\xrightarrow{\phi_Y}Y_A\xrightarrow{\lambda_A}X_A
       \xrightarrow{\phi_X^*}X'
\]
is \(n\lambda_{E\times F}\).
The last map is an isomorphism by the torus calculation, and the two
principal-polarization maps are isomorphisms. Hence \(\phi_Y\) has
absolute multiplier \(n\). Functoriality of the valuation pairing,
\[
 \langle\phi_Yy,x\rangle=\langle y,\phi_X^*x\rangle,
\]
now gives \(n\ell(A)=\ell(E\times F)\).
For the Tate parameter \(q_E\), the latter equals
\(-\log|q_E|_v=\log|j_E|_v\), by the expansion of \(j\)
\cite[Equation~(18)]{Tate}.
All absolute values are extended from the original field; this is
why the calculation is unchanged by the local extension.
\end{proof}

To use this reciprocal length in a height estimate on the curve,
we need to measure it by a boundary divisor on the moduli space.
The following comparison provides this link uniformly at finite places.

We specify the part of integral toroidal compactification theory used
in the proof. At full symplectic level \(m\ge3\), choose a smooth
projective toroidal compactification of the Siegel moduli scheme over
\(\Z[\zeta_m,1/m]\), with reduced boundary divisor \(D_m\).
The boundary is an effective relative Cartier divisor.

\begin{lemma}[Integral boundary comparison]\label{lem:boundary-comparison}
Fix the level \(m\) and compactification above. There are constants
\(c_m,C_m>0\), depending only on these choices, with the following
property. Let \(k_v/\Q_p\) be a finite extension with \(p\nmid m\),
and let \(A/k_v\) be a principally polarized abelian surface with
full level \(m\).
Write \(\lambda_{D_m,v}(A)\) for the model local height of the boundary.
\begin{enumerate}[label=(\roman*)]
\item If \(A\) has potential good reduction, then
\(\lambda_{D_m,v}(A)=0\).
\item If \(A\) has potential split semiabelian reduction of toric
rank one, then
\[
 c_m\ell(A)\le\lambda_{D_m,v}(A)\le C_m\ell(A).
\]
\end{enumerate}
Here \(\ell(A)\) is evaluated after an extension giving split
semiabelian reduction. With the absolute value extended from the
original field, both sides are unchanged by every further local
extension.
\end{lemma}

\begin{proof}
Here is the precise chart calculation behind these assertions.
The integral toroidal construction is developed in \cite{FC}; we use
\cite[Construction~6.3.1.1 and Theorem~6.4.1.1]{Lan}, with projectivity
supplied by \cite[Theorem~7.3.3.4]{Lan}.
By \cite[Theorem~6.4.1.1(2), (5)]{Lan}, the completed neighborhood of
a stratum is the prescribed torus embedding over an abelian-scheme
torsor over a finite \'etale cover of a lower-dimensional moduli
space. The boundary is the reduced relative Cartier normal-crossings
divisor of that embedding \cite[Theorem~6.4.1.1(3)]{Lan}.
Its strata pull back scheme-theoretically to the toroidal strata,
and their normal coordinates are the toric monomials
\cite[Proposition~6.3.1.6(1)]{Lan}.
After trivializing the torus torsor on a smooth cone chart, extend
the primitive ray generators to a lattice basis. In the resulting
coordinates \(q_1,\ldots,q_r\), its reduced-boundary ideal is
\((q_1\cdots q_r)\); this is the relative toric divisor of
\cite[Theorem~6.1.2.8(5)]{Lan}.

For a rank-one degeneration, the abelian quotient has good reduction.
Its moduli point therefore lies in the interior of this lower-dimensional
base. The extension coordinates belong to the abelian-scheme torsor;
after a local extension they extend over the valuation ring by
properness. This applies equally at ordinary and supersingular points
of that base. Thus there is no additional boundary coordinate from
the base or extension data.

The valuation form lies on a rank-one ray of the cone. If a rank-one
positive semidefinite form is a positive sum of such forms, each
summand vanishes on its kernel and lies on the same ray. Consequently
the ray is a face of every cone containing the form. In a smooth
cone chart only its primitive normal monomial, say \(q_c\), has
positive valuation; the other normal monomials are units. Locally
the reduced-boundary ideal is therefore \((q_c)\), up to a unit.
There is no vertical factor: this is the scheme-theoretic toric
boundary ideal in a relative normal-crossings chart over
\(\Z[\zeta_m,1/m]\), not a comparison only on the generic fiber.

Finally, \cite[Propositions~6.2.5.8(3), 6.2.5.11, and~6.3.1.6(4)]{Lan}
identify the tautological
ideal-valued toric pairing with the pairing of the semiabelian
degeneration. This identifies \(-\log|q_c|\) with \(\ell(A)\) times
the fixed rational factor comparing the primitive normal character
and the rank-one polarized cusp lattice. Those lattices and their
indices are part of the fixed level datum. They do not depend on
an isogeny presenting \(A\), or on the ramification of its field.
There are finitely many cusp types at fixed level, giving the stated
two fixed comparison constants. For good reduction the point is
interior and every boundary monomial is a unit. Computing these
values with the extended absolute value preserves them under every
local field extension.

The integral Siegel model here requires only that \(m\) be invertible,
including at residue characteristic two. This also follows from
the good-prime conditions in \cite[Definition~1.4.1.1]{Lan}:
for the self-dual symplectic datum over \(\Z\), the algebra is
\(\Q\), the order is \(\Z\), and there is no type-D or discriminant
exclusion. Later we will use levels \(3\) and \(4\), so every finite
prime is handled by an invertible level.
\end{proof}

\section{A uniform height estimate on a curve}\label{sec:height}

Fix a smooth curve \(S\) over a number field \(K_0\), carrying a
principally polarized abelian scheme of relative dimension two with
full level divisible by \(12\). Let \(\bar S\) be its smooth
projective completion. The two level-forgetting maps extend, by
properness, to morphisms
\[
 f_m:\bar S\longrightarrow\overline{\A}_{2,m}\qquad(m=3,4).
\]
Let \(\Delta\) be the reduced inverse image of the boundary.
Its support is the same at the two levels, because it is characterized
by nonzero potential toric rank. It may be empty.

\begin{proposition}\label{prop:height}
There are constants \(c,n_0>0\), depending only on this fixed family,
such that at every mixed point \(t\in S(\Qbar)\), with the notation
of Lemma~\ref{lem:gluing} and \(D=|\disc\End(F)|\),
\[
 n\ge n_0\quad\Longrightarrow\quad h(j_E)\le c\,n(1+\sqrt D).
\]
\end{proposition}

No Hodge-genericity hypothesis is needed for this height estimate.
When the boundary is nonempty, a height associated with it controls
\(h(j_E)\), with a logarithmic isogeny error. We will obtain a reverse
bound in which \(h(j_E)\) occurs with coefficient \(O(1/n)\), allowing
absorption for large \(n\). When the boundary is empty, the finite-place
estimates and compactness at infinity give the bound directly.

\subsection{Uniform local height comparisons}\label{subsec:model}

We first make explicit the uniformity in the height machine.
An \(M_{K_0}\)-bound means constants \(c_v\), zero at all but finitely
many finite places of \(K_0\), inherited at places above \(v\) by
extension of its absolute value. Their weighted sum is bounded
independently of any finite extension of \(K_0\).

If \(\Delta\ne\varnothing\), write
\(f_m^*D_m=\Delta+R_m\), where \(R_m\) is effective. Local heights of
effective divisors on a proper variety are bounded below.
Functoriality therefore gives
\begin{equation}\label{eq:local-comparison}
 \lambda_{\Delta,v}(t)
 \le\lambda_{D_m,v}(f_m(t))+O_v(1).
\end{equation}
If \(\Delta=\varnothing\), the generic pullback divisor is zero, so
instead
\begin{equation}\label{eq:zero-comparison}
 \lambda_{D_m,v}(f_m(t))=O_v(1).
\end{equation}
These errors are \(M_{K_0}\)-bounds. At finite places prime to \(m\),
we may use exactly the model local height described in
Lemma~\ref{lem:boundary-comparison}.

Here is a model justification that also handles the finitely many
bad primes of the curve and of its maps. Over
\(\mathcal O_{K_0}[1/m]\), take a proper model of \(\bar S\),
then the closure of the graph of \(f_m\) in its product with the
fixed integral toroidal compactification. If needed, pass to a
fixed further model on which the relevant divisors are Cartier. Models of the same generic divisor differ by
vertical divisors at finitely many primes. An effective vertical
Cartier divisor \(V\) is bounded above by a fixed multiple of the
full special fiber: locally a power of the uniformizer lies in the
ideal of \(V\), and quasi-compactness supplies a common power.
Taking positive and negative bounds for vertical discrepancies
shows that changing models introduces an error bounded in the
extended absolute value. This proves the asserted uniformity even
when the point and the ramification of its field vary.

Combining \eqref{eq:local-comparison} with
Lemmas~\ref{lem:boundary-comparison} and~\ref{lem:length},
and choosing \(m\in\{3,4\}\) prime to the residue characteristic,
gives, when \(\Delta\ne\varnothing\),
\begin{equation}\label{eq:finite-bound}
 \lambda_{\Delta,v}(t)
 \le c\,\frac{\log^+|j_E|_v}{n}+O_v(1),
\end{equation}
with fixed \(c\). If \(\Delta=\varnothing\), the two-sided comparison in
Lemma~\ref{lem:boundary-comparison} and \eqref{eq:zero-comparison} instead give
\begin{equation}\label{eq:finite-empty}
 \frac{\log^+|j_E|_v}{n}\le O_v(1).
\end{equation}
The bounds include the potential-good-reduction case, where the
boundary value and \(\log^+|j_E|_v\) both vanish.

\subsection{An archimedean boundary estimate}

At a complex point, let \(\rho(A)\) be the length of the shortest
nonzero period of \(A\) for its principal polarization.

\begin{lemma}\label{lem:arch-boundary}
If \(\Delta\ne\varnothing\), then at each archimedean place of \(K_0\),
\[
 \lambda_\Delta(t)\ll 1+\rho(A_t)^{-2}.
\]
\end{lemma}

\begin{proof}
Work near one point of \(\Delta\), with disk parameter \(z\).
After a fixed ramified base change, monodromy on \(H_1\) is unipotent
and has logarithm \(N\) with \(N^2=0\); this is the monodromy
description for a semistable degeneration of abelian varieties
\cite{SGA7}.
The image of \(N\) is isotropic. Choose a rational Lagrangian
containing it and contained in \(\ker N\), and adapt an integral
symplectic basis to its primitive lattice. On the universal cover
of the punctured disk, the resulting period matrix \(Z\) has monodromy
\[
 Z\longmapsto Z+M,\qquad M\in\operatorname{Sym}_2(\Z).
\]

For every \(u\in\Z^2\), the holomorphic function
\(q_u=\exp(2\pi i\,u^tZu)\) is single-valued and bounded in absolute
value by one. It extends across \(z=0\); its order there is
\(u^tMu\), by the monodromy of its logarithm.
Thus \(M\) is positive semidefinite. Writing
\(q_u=z^{u^tMu}\) times a holomorphic unit, and using the finitely
many vectors \(e_i,e_i+e_j\), shows that
\begin{equation}\label{eq:period-boundary}
 Y:=\Im Z=\frac{-\log|z|}{2\pi}M+R(z),
\end{equation}
where \(R\) extends harmonically across zero.

On \(\ker M\), the form \(R(0)\) is positive definite. For any
nonzero real \(v\) in that kernel, \(v^tR(z)v=v^tYv>0\) on the
punctured disk. If its value at zero were zero, the minimum
principle would force this harmonic function to vanish identically.
If \(M=0\), the same argument and the holomorphic extension just
constructed would extend \(Z\) to the interior of Siegel space.
Uniqueness of extension to the separated toroidal compactification
would contradict the choice of a boundary point. Hence \(M\ne0\).

Choose a nonzero integral \(w\in\im M\), possible because \(M\) is
integral. On the orthogonal decomposition
\(\im M\oplus\ker M\), the first block of \(Y\) grows as a positive
definite matrix times \(-\log|z|\), the off-diagonal blocks are
bounded, and the kernel block stays positive definite.
The Schur-complement formula gives
\[
 w^tY^{-1}w=O\bigl((-\log|z|)^{-1}\bigr).
\]
This is the squared norm of the period \(w\) in
\(\Z^2+Z\Z^2\). Therefore
\(\rho(A_t)^2\ll(-\log|z|)^{-1}\), and
\(-\log|z|\ll\rho(A_t)^{-2}\).
The boundary local height is a fixed multiple of \(-\log|z|\) up to
a bounded term, also after the fixed ramification.
Finitely many such disks, followed by compactness on their
complement, prove the lemma. The argument includes both ranks
one and two for \(M\).
\end{proof}

\begin{lemma}\label{lem:rho}
At a mixed point, choose reduced elliptic periods at a complex
embedding and write \(y_E=\Im\tau_E\), \(y_F=\Im\tau_F\). Then
\[
 \rho(A)^{-2}\le ny_F+\frac{y_E}{n},
 \qquad
 \rho(A)^2\le\frac n{y_E}.
\]
\end{lemma}

\begin{proof}
Identify Lie spaces using \(E\times F\to A\).
The period lattice of \(A\) is an overlattice of
\(\Lambda_E\oplus\Lambda_F\) contained in its \(1/n\) multiple, and
its Hermitian form is \(n\) times the product form.
Its intersection with \(\Lie E\) is exactly \(\Lambda_E\),
because the graph kernel has trivial intersection with
\(E\times\{0\}\).

A period with nonzero \(F\)-projection has squared norm at least
\(n\cdot n^{-2}/y_F=1/(ny_F)\).
A nonzero period with zero \(F\)-projection has squared norm at
least \(n/y_E\), and the first reduced period of \(E\) attains this
value. The two assertions follow.
\end{proof}

\subsection{Proof of Proposition~\ref{prop:height}}

Suppose first that \(\Delta=\varnothing\).
At every embedding of \(K_0\), the image of \(\bar S\) is compact
in the interior moduli space. The shortest period length has a
positive lower bound on this image. Lemma~\ref{lem:rho} therefore
gives \(y_E\ll n\) at every embedding of a field of definition of
the point. Equations~\eqref{eq:j-y} and \eqref{eq:finite-empty},
summed with the normalized local weights, show
\[
 h(j_E)\ll n.
\]
This proves the proposition in this case.

Now let \(\Delta\ne\varnothing\). Choose \(h_\Delta\) shifted to be
nonnegative. It is an ample height because \(\Delta\) is a nonzero
effective divisor on a projective curve.
Lemmas~\ref{lem:arch-boundary} and \ref{lem:rho}, together with
\eqref{eq:b-index}, give at the archimedean places
\[
 \lambda_\Delta(t)\ll
 1+n\sqrt D+\frac{y_E}{n}.
\]
Use \eqref{eq:j-y} for \(y_E\), add \eqref{eq:finite-bound}, and sum.
The \(M_{K_0}\)-bounds give a constant independent of the field.
We obtain
\begin{equation}\label{eq:height-upper}
 h_\Delta(t)\le B\left(1+n\sqrt D+\frac{h(j_E)}n\right)
\end{equation}
for a fixed \(B>0\).

On the other hand, \eqref{eq:faltings-isogeny} gives
\(\hF(E)+\hF(F)\le\hF(A)+\log n\).
The lower bound for elliptic Faltings height,
\eqref{eq:elliptic-height}, and Lemma~\ref{lem:family-height} imply
\begin{equation}\label{eq:height-lower}
 h(j_E)\le A_0\bigl(1+h_\Delta(t)+\log n\bigr)
\end{equation}
for fixed \(A_0>0\). Combining \eqref{eq:height-upper} and
\eqref{eq:height-lower}, the coefficient of \(h(j_E)\) on the
right is \(A_0B/n\).
For \(n\ge2A_0B\) it can be absorbed.
Since \(\log n\ll n\), this gives
\(h(j_E)\ll n(1+\sqrt D)\), as claimed.
\qed

\section{An essential period on a threefold}\label{sec:threefold}

We now extract a stronger power of the gluing index by repeating the
non-CM factor. This lets one period vector contain both independent
periods of \(E\), so no nonzero integer row on \(E^2\) can annihilate
its \(E\)-coordinates. Repeated elliptic factors and polarizations
adapted to reduced period bases are used in
\cite[\S\S7.2--7.3]{GR}; here the gluing data produce a threefold
isogenous to \(E^2\times F\).
All constructions in this section are over a number
field \(k\) containing the definitions of \(E,F,\theta\), and all
endomorphisms of \(F\). Put \(d=[k:\Q]\), and fix one complex
embedding. Take a reduced basis \(e_1,e_2\) of \(\Lambda_E\), and
a first reduced period \(f\) of \(F\). Use \(b,y_F\) from
Lemma~\ref{lem:cm}, and put \(y=y_E\).

The two squared period lengths of \(E\) have sizes \(1/y\) and
\(y\). Weighting the first copy of \(E\) by approximately \(y^2\)
makes their contributions comparable. Weighting \(F\) by approximately
\(y\) then makes the cube root of the product polarization degree
grow at the same rate as the squared norm. The remaining gain is
the power \(n^{4/3}\) contributed by the quotient.

\begin{proposition}\label{prop:threefold}
There is a polarized abelian threefold \((B,L)/k\), isogenous to
\(E^2\times F\), such that
\begin{equation}\label{eq:threefold-estimate}
 n^{4/3}\ll
 (b^2+1)d\,\max\{1,\hF(B),\log\deg_L B\},
\end{equation}
and
\begin{equation}\label{eq:B-height-degree}
 \hF(B)\le2\hF(E)+\hF(F)+\log n,\qquad
 \deg_LB=6n^4uv,
 \quad u=\ceil{y^2},\quad v=\ceil y.
\end{equation}
\end{proposition}

\begin{proof}
We use the index \(b\) to realize both points
\(\theta^{-1}(be_i/n)\) as endomorphism images of the common
torsion point \(f/n\). For \(i=1,2\), choose \(f_i\in\Lambda_F\)
so that the class of \(f_i/n\) is \(\theta^{-1}(e_i/n)\).
Since \(b\Lambda_F\subset\End(F)f\), there is \(\beta_i\in\End(F)\)
with
\begin{equation}\label{eq:beta}
 \beta_if=bf_i.
\end{equation}
There is no need for \(\beta_i\) to be invertible on \(F[n]\).

Define
\[
 \Gamma=\{(\theta\beta_1(p),\theta\beta_2(p),p):p\in F[n]\},
 \qquad B=(E^2\times F)/\Gamma.
\]
All maps in this expression are defined over \(k\).
Indeed, the analytically selected \(\beta_i\) belong to \(\End(F)\),
all of which is defined over \(k\). The last-coordinate projection
is inverse to the displayed graph embedding as a group-scheme map.
It therefore identifies \(\Gamma\) with \(F[n]\), without any
coprimality requirement on \(n,b,\beta_1,\beta_2\); hence
the quotient map \(\pi:E^2\times F\to B\) has degree \(n^2\).
The quotient and its map are defined over \(k\). Since \(\Gamma\)
is killed by \(n\), the quotient property factors multiplication
by \(n\), over the same field, as
\[
 E^2\times F\xrightarrow{\pi}B\xrightarrow{\psi}E^2\times F,
 \qquad\psi\pi=[n].
\]
Since \(\deg[n]=n^6\), one has \(\deg\psi=n^4\).

Give the product the line bundle
\[
 M=\mathcal O_E(u[0])\boxtimes
    \mathcal O_E([0])\boxtimes\mathcal O_F(v[0]),
 \qquad L=\psi^*M.
\]
These ample line bundles are defined over \(k\).
Although the integers \(u,v\) were chosen using one embedding,
they require no extension of the field. Since
\(\deg_M(E^2\times F)=6uv\), the degree assertion follows.
The height assertion follows from \(\deg\pi=n^2\) and
\eqref{eq:faltings-isogeny}.

In the product Lie space, the vector
\[
 \omega=\frac{(be_1,be_2,f)}n
\]
is a period of \(B\). Indeed \eqref{eq:beta} gives
\(\theta\beta_i(f/n)=be_i/n\), precisely the membership condition
for the quotient lattice.

This period is essential with respect to every geometric abelian
subvariety. If it lay in the Lie algebra of a proper abelian
subvariety of \(B_\C\), the identity component of its inverse image
under \(\pi\) would be a proper abelian subvariety of
\(E^2\times F\) whose Lie algebra contains the displayed vector. By semisimplicity up to isogeny, such a
subvariety is annihilated by a nonzero homomorphism to \(E\) or \(F\).
A homomorphism to \(E\) is an integer row \((a_1,a_2,0)\), since
\(\End(E)=\Z\) and \(\Hom(F,E)=0\); its value on the vector cannot
vanish, because \(e_1,e_2\) are independent over \(\Z\).
A homomorphism to \(F\) acts by a nonzero endomorphism on \(f\),
and again cannot vanish. A rational homomorphism may be made
integral by clearing denominators, which does not affect this
argument.

The identity \(\pi^*L=[n]^*M\) gives the norm
\[
 \|\omega\|_L^2
 =b^2u\|e_1\|^2+b^2\|e_2\|^2+v\|f\|^2
 \ll(b^2+1)y.
\]
Here \eqref{eq:elliptic-norms} bounds \(u/y\) and
\(|\tau_E|^2/y\) by \(O(y)\), and \(v/y_F\ll y\).
The factor \(n^2\) in the pulled-back metric has cancelled the
denominator \(n^2\) in the squared norm.
At the same time, \((uv)^{1/3}\ge y\), so
\[
 \frac{(\deg_LB)^{1/3}}{\|\omega\|_L^2}
 \gg\frac{n^{4/3}}{b^2+1}.
\]
Apply \eqref{eq:period-theorem} to this essential period.
\end{proof}

\begin{remark}
The large endomorphisms \(\beta_i\) do not appear in the bound.
Their role is to put a specified vector into the quotient lattice,
not to bound the degree of a homomorphism. The projection of
\(\Gamma\) to \(F[n]\) fixes its order even if either \(\beta_i\)
has a large kernel on \(F[n]\).
\end{remark}

\section{The Galois orbit bound}\label{sec:orbits}

We apply the preceding estimates to the curve of
Theorem~\ref{thm:orbits}. Pass to a component of a full-level cover
with level divisible by \(12\), and normalize it. After choosing a
number field \(K_0\), enlarged once and for all, this gives a smooth
curve \(S/K_0\) with a family as in Section~\ref{sec:height}, and a
map \(S\to C\) of fixed bounded degree over a dense open.
Further fixed covers or restrictions used to describe the family
have the same property. Only finitely many points of \(C\) are
omitted. This passage is permitted for a singular \(C\) as well.

If \(s\) belongs to that dense open, choose a lift \(t\) with
\[
 [K_0(t):\Q]\ll [K(s):K].
\]
Such a lift exists because the finite fiber has bounded degree,
and \(K_0/K\) is fixed. By Lemma~\ref{lem:gluing}, \(E,F,\theta\)
are defined over the field of the fiber. After an extension of
degree at most two, all CM endomorphisms are defined too.
For the resulting field \(k\), put \(d=[k:\Q]\); then
\begin{equation}\label{eq:field-cost}
 d\ll [K(s):K].
\end{equation}
It remains to bound \(\max\{n^2,D\}\) by a fixed power of \(d\).

For \(n\ge n_0\), Lemma~\ref{lem:cm} and
Proposition~\ref{prop:height} give
\begin{equation}\label{eq:coarse-height}
 D\ll d^4,\qquad b^2+1\ll d^4,\qquad
 h(j_E)\ll n d^2,\qquad h(j_F)\ll d^2.
\end{equation}
Let \((B,L)\) be the threefold from Proposition~\ref{prop:threefold}.
Its height satisfies
\[
 \hF(B)\le2\hF(E)+\hF(F)+\log n\ll nd^2.
\]
To bound its polarization degree, use \eqref{eq:j-y} at the chosen
embedding. A single archimedean contribution to absolute height
is at most \(d\) times the total, so
\[
 y\ll1+d\,h(j_E)\ll nd^3.
\]
It follows that
\(\log(6n^4uv)\ll1+\log n+\log d\ll nd^2\).
Thus \eqref{eq:threefold-estimate} and \eqref{eq:coarse-height} give
\[
 n^{4/3}\ll d^4\cdot d\cdot nd^2=nd^7.
\]
Dividing by \(n\) and cubing proves
\begin{equation}\label{eq:coarse-orbit}
 n\ll d^{21},\qquad
 \max\{n^2,D\}\ll d^{42}.
\end{equation}
For the remaining \(n<n_0\), the same final bound follows from
\(D\ll d^4\) by increasing the constant.

By Lemma~\ref{lem:gluing}, \(\cplx(s)=\max\{n^2,D\}\).
Equations~\eqref{eq:field-cost} and \eqref{eq:coarse-orbit} therefore
imply
\[
 [K(s):K]\gg\cplx(s)^{1/42}.
\]
The degree on the left is the cardinality of the Galois orbit.
The finitely many omitted mixed points can be included by reducing
the positive constant. This proves Theorem~\ref{thm:orbits}.
\qed

\section{Finiteness on the original curve}\label{sec:finiteness}

We recall exactly the unlikely-intersection implication being used.
Daw--Orr \cite[Theorem~1.2]{DO} proves that a Hodge-generic curve in
\(\A_2\) has finitely many mixed \(E\times\CM\) points if their
Galois orbits satisfy a power lower bound in the relevant special
curve complexity. Their alternative complexity is the maximum of
the minimum product-isogeny degree and the endomorphism-ring
discriminant of the CM factor in a product attaining that degree;
it is polynomially comparable with their original complexity
\cite[\S3, especially Lemmas~3.1 and~3.6]{DO}.
It is exactly \(\cplx(s)\) by Lemma~\ref{lem:gluing}.
Theorem~\ref{thm:orbits} supplies the required bound.
No boundary hypothesis is present in this implication.\footnote{We use
the corrected version of \cite{DO}. Its Lemma~2.3 verifies the
Cartan-stability condition required in the corrected reduction-theory
input \cite{OS}.}

The mixed points occurring in the Daw--Orr formulation are algebraic
in our setting. Indeed, a mixed complex point
lies on an \(E\times\CM\) special curve defined over \(\Qbar\):
start with the diagonal product locus with a fixed CM elliptic
coordinate and use the polarized isogeny of
Lemma~\ref{lem:gluing} to pass through a Hecke correspondence.
Such a special curve meets \(C\) in a zero-dimensional set, because
\(C\) is Hodge generic. These intersection points are algebraic.
The elliptic subvarieties and their homomorphisms are then available
over \(\Qbar\), by invariance of homomorphisms of abelian varieties
under algebraically closed extension. Thus the complex and algebraic
mixed loci required here coincide.

The same orbit bound supplies the finitely-generated-field version
of Daw--Orr's hypothesis for this curve over \(\Qbar\).
Given a finitely generated field of definition \(L\subset\C\), let
\(K\) be a number field of definition of \(C\), set \(M=LK\), and
put \(K'=M\cap\Qbar\). Then \(K'\) is a number field containing
\(K\), and \(M/K'\) is regular. For an algebraic point \(s\),
linear disjointness shows that its \(\operatorname{Aut}(\C/M)\)-orbit
has cardinality \([K'(s):K']\). Theorem~\ref{thm:orbits} applies
over \(K'\). Since \(\operatorname{Aut}(\C/M)\) is a subgroup of
\(\operatorname{Aut}(\C/L)\), the required lower bound holds for
the latter orbit too.

We conclude that \(C\) has finitely many mixed points. If both
elliptic factors have CM, their product, and hence every isogenous
abelian surface, is of CM type. Its moduli point is special.
Andr\'e--Oort for \(\A_2\) gives only finitely many special points
on a Hodge-generic curve. Together these two conclusions prove
Theorem~\ref{thm:main} on \(C\) itself.
\qed

\begin{thebibliography}{99}

\bibitem{Chen2007}
Y.-G. Chen,
\emph{On the Siegel--Tatuzawa--Hoffstein theorem},
Acta Arith. \textbf{130} (2007), 361--367.
\href{https://www.impan.pl/shop/en/publication/transaction/download/product/82398}
{Publisher's PDF}.

\bibitem{Cox}
D. A. Cox,
\emph{Primes of the Form \(x^2+ny^2\): Fermat, Class Field Theory,
and Complex Multiplication},
second ed., John Wiley \& Sons, Hoboken, NJ, 2013.
\href{https://doi.org/10.1002/9781118400722}{doi:10.1002/9781118400722}.

\bibitem{DO}
C. Daw and M. Orr,
\emph{Unlikely intersections with \(E\times\CM\) curves in \(\A_2\)},
Ann. Sc. Norm. Super. Pisa Cl. Sci. (5) \textbf{22} (2021), 1705--1745.
\href{https://doi.org/10.2422/2036-2145.202006_014}{doi:10.2422/2036-2145.202006\_014}.
We use the corrected preprint,
\href{https://arxiv.org/abs/1902.10483v4}{arXiv:1902.10483v4}
(28 May 2021).

\bibitem{DOMultiplicative}
C. Daw and M. Orr,
\emph{The large Galois orbits conjecture under multiplicative degeneration},
\href{https://arxiv.org/abs/2306.13463v2}{arXiv:2306.13463v2},
revised 15 May 2025.

\bibitem{FC}
G. Faltings and C.-L. Chai,
\emph{Degeneration of Abelian Varieties},
Ergebnisse der Mathematik und ihrer Grenzgebiete (3), vol.~22,
Springer-Verlag, Berlin, 1990.

\bibitem{GR}
\'E. Gaudron and G. R\'emond,
\emph{Th\'eor\`eme des p\'eriodes et degr\'es minimaux d'isog\'enies},
Comment. Math. Helv. \textbf{89} (2014), 343--403.
\href{https://doi.org/10.4171/CMH/322}{doi:10.4171/CMH/322}.

\bibitem{SGA7}
A. Grothendieck,
\emph{Mod\`eles de N\'eron et monodromie}, with an appendix by M. Raynaud,
Expos\'e IX in \emph{Groupes de monodromie en g\'eom\'etrie alg\'ebrique}
(SGA~7~I), Lecture Notes in Mathematics, vol.~288,
Springer, Berlin, 1972, 313--523.
\href{https://doi.org/10.1007/BFb0068694}{doi:10.1007/BFb0068694}.

\bibitem{Lan}
K.-W. Lan,
\emph{Arithmetic Compactifications of PEL-Type Shimura Varieties},
London Mathematical Society Monographs Series, vol.~36,
Princeton University Press, Princeton, 2013.
Numbered references use the
\href{https://www.kwlan.org/articles/cpt-PEL-type-thesis-revision.pdf}
{revised dissertation version} (14 March 2021).

\bibitem{Lobrich}
S. L\"obrich,
\emph{A gap in the spectrum of the Faltings height},
J. Th\'eor. Nombres Bordeaux \textbf{29} (2017), 289--305.
\href{https://doi.org/10.5802/jtnb.980}{doi:10.5802/jtnb.980}.

\bibitem{MW}
D. W. Masser and G. W\"ustholz,
\emph{Periods and minimal abelian subvarieties},
Ann. of Math. (2) \textbf{137} (1993), 407--458.
\href{https://doi.org/10.2307/2946542}{doi:10.2307/2946542}.

\bibitem{OpenAIQuasiRH2026}
\begingroup\raggedright
OpenAI,
\emph{The Quasi-Riemann Hypothesis: A Zero-Free Half-Plane
\(\mathrm{Re}(s)>7/8\)},
OpenAI Math Release preprint
\href{https://github.com/openai/math/blob/main/preprints/The-Quasi-Riemann-Hypothesis-September-30-2026/paper.pdf}{OAI:The-Quasi-Riemann-Hypothesis-September-30-2026},
2026.
\par\endgroup

\bibitem{OS}
M. Orr and C. Schnell,
\emph{Correction to the article Height bounds and the Siegel property},
Algebra Number Theory \textbf{17} (2023), 1231--1237.
\href{https://doi.org/10.2140/ant.2023.17.1231}{doi:10.2140/ant.2023.17.1231}.

\bibitem{Papas}
G. Papas,
\emph{On the \(v\)-adic values of \(G\)-functions I: Splittings in \(\A_2\)},
\href{https://arxiv.org/abs/2510.11806v1}{arXiv:2510.11806v1},
13 October 2025.

\bibitem{Pazuki}
F. Pazuki,
\emph{Theta height and Faltings height},
Bull. Soc. Math. France \textbf{140} (2012), 19--49.
\href{https://doi.org/10.24033/bsmf.2623}{doi:10.24033/bsmf.2623}.

\bibitem{PT}
J. Pila and J. Tsimerman,
\emph{The Andr\'e--Oort conjecture for the moduli space of abelian surfaces},
Compos. Math. \textbf{149} (2013), 204--216.
\href{https://doi.org/10.1112/S0010437X12000589}
{doi:10.1112/S0010437X12000589}.

\bibitem{PW}
J. Pila and A. J. Wilkie,
\emph{The rational points of a definable set},
Duke Math. J. \textbf{133} (2006), 591--616.
\href{https://doi.org/10.1215/S0012-7094-06-13336-7}
{doi:10.1215/S0012-7094-06-13336-7}.

\bibitem{PZ}
J. Pila and U. Zannier,
\emph{Rational points in periodic analytic sets and the Manin--Mumford conjecture},
Rend. Lincei Mat. Appl. (9) \textbf{19} (2008), 149--162.
\href{https://doi.org/10.4171/RLM/514}{doi:10.4171/RLM/514}.

\bibitem{Pink}
R. Pink,
\emph{A common generalization of the conjectures of Andr\'e--Oort,
Manin--Mumford, and Mordell--Lang}, preprint, April 17, 2005.
\href{https://people.math.ethz.ch/~pink/ftp/AOMMML.pdf}{Author's manuscript}.

\bibitem{ST}
J.-P. Serre and J. Tate,
\emph{Good reduction of abelian varieties},
Ann. of Math. (2) \textbf{88} (1968), 492--517.
\href{https://doi.org/10.2307/1970722}{doi:10.2307/1970722}.

\bibitem{Siegel}
C. L. Siegel,
\emph{\"Uber die Classenzahl quadratischer Zahlk\"orper},
Acta Arith. \textbf{1} (1935), 83--86.
\href{https://doi.org/10.4064/aa-1-1-83-86}{doi:10.4064/aa-1-1-83-86}.

\bibitem{Silverman}
J. H. Silverman,
\emph{Heights and elliptic curves}, in \emph{Arithmetic Geometry}
(G. Cornell and J. H. Silverman, eds.), Springer, New York, 1986,
253--265.
\href{https://doi.org/10.1007/978-1-4613-8655-1_10}
{doi:10.1007/978-1-4613-8655-1\_10}.

\bibitem{Tate}
J. Tate,
\emph{A review of non-Archimedean elliptic functions}, in
\emph{Elliptic Curves, Modular Forms, \& Fermat's Last Theorem}
(Hong Kong, 1993), Series in Number Theory, vol.~I,
International Press, Cambridge, MA, 1995, 162--184.
\href{https://people.math.harvard.edu/~ctm/home/text/others/tate/tate-unif/tate-unif.pdf}
{Author's reprint}.

\bibitem{Tsimerman}
J. Tsimerman,
\emph{The Andr\'e--Oort conjecture for \(\A_g\)},
Ann. of Math. (2) \textbf{187} (2018), 379--390.
\href{https://doi.org/10.4007/annals.2018.187.2.2}
{doi:10.4007/annals.2018.187.2.2}.

\bibitem{Zilber}
B. Zilber,
\emph{Exponential sums equations and the Schanuel conjecture},
J. London Math. Soc. (2) \textbf{65} (2002), 27--44.
\href{https://doi.org/10.1112/S0024610701002861}
{doi:10.1112/S0024610701002861}.

\end{thebibliography}
\end{document}
